% TOP_PROBLEM: {"id":30007565,"problem_number":"LOCAL-30007565","title":"Strategic adaptation after regime changes","category_id":21,"category":{"id":21,"name":"economics","display_name":"Economics","description":"Open economic theory statements, published research agendas, and proposed scoped specifications; question provenance and historical priority are qualified per record.","slug":"economics","order_index":21,"created_at":"2026-10-08T00:00:00Z"},"status":"research_question","external_url":null,"legacy_ids":[],"published":false,"statement":"Distinguish parameter learning from model revision after a randomized institutional payoff change, and identify effects on revision times and finite-horizon decision losses.","economics":{"original_id":54,"field":"Game theory and strategic interaction","kind":"identification","formalization_basis":"adapted_research_question","source_status":"research_agenda","priority_status":"earliest_found","priority_note":"This is the earliest source in which the relevant agenda was verified during this audit; absolute priority has not been established.","earliest_verified_reference":{"authors":["Drew Fudenberg","David K. Levine"],"title":"Whither Game Theory? Towards a Theory of Learning in Games","year":2016,"url":"https://economics.mit.edu/sites/default/files/2022-09/Whither%20Game%20Theory%20Towards%20a%20Theory%20of%20Learning%20in%20Games.pdf","doi":"10.1257/jep.30.4.151","locator":"Introduction, p.153; conclusion, pp.165–166","evidence_summary":"The essay explicitly calls for more knowledge about learning speed, equilibrium prediction, and tractable models that allow reevaluation."},"current_reference":{"authors":["Drew Fudenberg","David K. Levine"],"title":"Whither Game Theory? Towards a Theory of Learning in Games","year":2016,"url":"https://economics.mit.edu/sites/default/files/2022-09/Whither%20Game%20Theory%20Towards%20a%20Theory%20of%20Learning%20in%20Games.pdf","doi":"10.1257/jep.30.4.151","locator":"Introduction, p.153; conclusion, pp.165–166","evidence_summary":"The essay explicitly calls for more knowledge about learning speed, equilibrium prediction, and tractable models that allow reevaluation."},"formalization_note":"This is human model-reassessment identification. It does not duplicate uniform epsilon-equilibrium or stochastic-game decision problems."},"statusQualification":"Published question or research agenda verified at the cited locator. The mathematical specification is adapted for this catalogue and includes additional assumptions; source publication does not establish first-ever priority or certify this exact model remains unsolved. This is human model-reassessment identification. It does not duplicate uniform epsilon-equilibrium or stochastic-game decision problems.","statusReviewedAt":"2026-10-08"}
% FIRST_POSTED: 2026-10-08
% ENTRY_KIND: statement_only
% !TeX program = lualatex
\documentclass[11pt]{article}
\usepackage{fontspec}
\usepackage[a4paper,margin=25mm]{geometry}
\usepackage{amsmath,amssymb,mathtools}
\usepackage{xurl}
\usepackage[hidelinks]{hyperref}
\newcommand{\sref}[2]{\hyperlink{source:#1:#2}{[S#2]}}
\setlength{\emergencystretch}{3em}
\begin{document}
% problemId:30007565
\section*{Strategic adaptation after regime changes}\label{30007565}
\subsection{Definitions and mathematical statement}
Let human players repeatedly interact in a finite game with known action sets. At an experimentally randomized date \(\tau\), an institution changes the payoff matrix from \(u^0\) to \(u^1\). Players may receive an announcement or learn the change only through realized outcomes; this disclosure \(Z\) is randomized by interacting group. A candidate learning model maintains latent state \(M_{it}\in\{\text{old model},\text{revised model}\}\) and parameter estimates conditional on that state. Observe actions, realized payoffs, and incentivized predictive beliefs \(b_{it}\). Define revision time \(R_i\) by a prespecified belief-consistency criterion, rather than by an unobservable analyst judgment. The principal estimands are
\[
 \Pr(R_i-\tau\leq h\mid do(Z))
 \quad\text{and}\quad
 E\left[\sum_{t=\tau}^{\tau+h}(u_i^{\mathrm{best}}(t)-u_i(a_t))\mid do(Z)\right],
\]
where \(h\) is a fixed horizon and \(u_i^{\mathrm{best}}(t)=\max_{a_i}u_i^1(a_i,a_{-i,t})\). Determine whether observations distinguish changing beliefs within one model from switching the model itself, and derive identified sets when they cannot. Estimate disclosure effects and validate revision triggers on new changes. Regret here is a finite experimental outcome, not an equilibrium-existence target. Group randomization and explicit signal measurement assumptions are necessary because opponents also adapt after the change.

\par\textbf{Formulation and scope.} This is human model-reassessment identification. It does not duplicate uniform epsilon-equilibrium or stochastic-game decision problems.

\par\textbf{Open-question provenance.} An explicit question or research agenda was verified at the source locator. This does not by itself certify that every later version remains unresolved \sref{30007565}{1}.

\par\textbf{Historical priority.} This is the earliest source in which the relevant agenda was verified during this audit; absolute priority has not been established.

\subsection{Short English statement}
Distinguish parameter learning from model revision after a randomized institutional payoff change, and identify effects on revision times and finite-horizon decision losses.

\subsection{Sources}
\begin{itemize}\raggedright
\item[\textnormal{[S1]}]\hypertarget{source:30007565:1}{} Drew Fudenberg, David K. Levine. 2016. Whither Game Theory? Towards a Theory of Learning in Games. Introduction, p.153; conclusion, pp.165–166.
\url{https://economics.mit.edu/sites/default/files/2022-09/Whither%20Game%20Theory%20Towards%20a%20Theory%20of%20Learning%20in%20Games.pdf}.
DOI: 10.1257/jep.30.4.151.
{\small\textit{Earliest verified question/agenda reference; historical priority qualified below. The essay explicitly calls for more knowledge about learning speed, equilibrium prediction, and tractable models that allow reevaluation.}}
\end{itemize}
\end{document}
